% =============================================================================
% ARIS standardized research review report template
% Engine: xelatex (or latexmk -xelatex)
% Encoding: UTF-8
% =============================================================================

\documentclass[11pt,a4paper]{article}

% -----------------------------------------------------------------------------
% Metadata to fill for each report
% -----------------------------------------------------------------------------
\newcommand{\ReportTitle}{审稿报告}
\newcommand{\PaperTitle}{TODO: Paper title}
\newcommand{\PaperSubtitle}{TODO: English title or short subtitle}
\newcommand{\PaperShortName}{TODO: short name}
\newcommand{\VenueName}{TODO: venue or journal}
\newcommand{\ManuscriptId}{TODO: manuscript id}
\newcommand{\ReviewerName}{匿名审稿人 (ARIS Auto-Review Pipeline)}
\newcommand{\ReportDate}{\today}
\newcommand{\Recommendation}{TODO: Accept / Minor Revision / Major Revision / Reject}
\newcommand{\OverallScore}{TODO: score}
\newcommand{\ReviewConfidence}{TODO: High / Medium / Low}
\newcommand{\ReviewScope}{TODO: manuscript, supplementary files, code/data, literature search, external review}

% -----------------------------------------------------------------------------
% Packages and typography
% -----------------------------------------------------------------------------
\usepackage[a4paper,left=2.4cm,right=2.4cm,top=2.6cm,bottom=2.6cm]{geometry}
\usepackage{amsmath,amssymb,amsfonts,amsthm}
\usepackage{bm}
\usepackage{graphicx}
\usepackage{booktabs}
\usepackage{array}
\usepackage{longtable}
\usepackage{multirow}
\usepackage{tabularx}
\usepackage{enumitem}
\usepackage{xcolor}
\usepackage{tikz}
\usetikzlibrary{shapes,arrows,positioning,calc}
\usepackage[most]{tcolorbox}
\usepackage{fancyhdr}
\usepackage{titlesec}
\usepackage{caption}
\usepackage{microtype}
\usepackage{adjustbox}
\usepackage{ragged2e}
\usepackage[UTF8]{ctex}
\usepackage{hyperref}

\hypersetup{
  colorlinks=true,
  linkcolor=blue!60!black,
  citecolor=blue!60!black,
  urlcolor=blue!60!black,
  pdfauthor={\ReviewerName},
  pdftitle={\PaperShortName\space 审稿报告}
}

\pagestyle{fancy}
\fancyhf{}
\fancyhead[L]{\small \PaperShortName\space 审稿报告}
\fancyhead[R]{\small \thepage}
\fancyfoot[C]{\small ARIS Auto-Review Pipeline \quad|\quad \ManuscriptId}
\renewcommand{\headrulewidth}{0.4pt}
\renewcommand{\footrulewidth}{0.4pt}

\titleformat{\section}{\Large\bfseries\color{blue!50!black}}{\thesection.}{0.6em}{}
\titleformat{\subsection}{\large\bfseries\color{blue!40!black}}{\thesubsection.}{0.5em}{}
\titleformat{\subsubsection}{\normalsize\bfseries\color{blue!30!black}}{\thesubsubsection.}{0.5em}{}

\setlength{\parskip}{0.4em}
\setlength{\parindent}{2em}
\linespread{1.22}
\setlist[itemize]{leftmargin=2em,itemsep=0.2em,topsep=0.2em}
\setlist[enumerate]{leftmargin=2em,itemsep=0.2em,topsep=0.2em}

% -----------------------------------------------------------------------------
% Reusable review boxes
% -----------------------------------------------------------------------------
\newtcolorbox{summarybox}[1][]{
  colback=blue!5!white,colframe=blue!55!black,
  fonttitle=\bfseries,title={#1},
  arc=1mm,boxrule=0.45pt,breakable
}

\newtcolorbox{positivebox}[1][]{
  colback=green!6!white,colframe=green!50!black,
  fonttitle=\bfseries,title={#1},
  arc=1mm,boxrule=0.45pt,breakable
}

\newtcolorbox{criticalbox}[1][]{
  colback=red!6!white,colframe=red!60!black,
  fonttitle=\bfseries,title={#1},
  arc=1mm,boxrule=0.45pt,breakable
}

\newtcolorbox{evidencebox}[1][]{
  colback=yellow!8!white,colframe=yellow!60!black,
  fonttitle=\bfseries,title={#1},
  arc=1mm,boxrule=0.45pt,breakable
}

\newtcolorbox{innovationbox}[1][]{
  colback=cyan!6!white,colframe=cyan!55!black,
  fonttitle=\bfseries,title={#1},
  arc=1mm,boxrule=0.45pt,breakable
}

\newcommand{\IssueHeading}[3]{\noindent\textbf{#1. #2} \hfill \textbf{Severity: #3}\par}
\newcommand{\Evidence}{\noindent\textbf{Evidence:} }
\newcommand{\Impact}{\noindent\textbf{Impact:} }
\newcommand{\RequiredAction}{\noindent\textbf{Required action:} }
\newcommand{\Confidence}{\noindent\textbf{Confidence:} }

% -----------------------------------------------------------------------------
\begin{document}

\begin{titlepage}
\vspace*{-1.0cm}
\begin{tcolorbox}[colback=blue!5,colframe=blue!60!black,arc=2mm,boxrule=0.6pt]
\centering
{\LARGE\bfseries \ReportTitle\par}
\vspace{0.5em}
{\large \PaperTitle\par}
\vspace{0.3em}
{\normalsize \PaperSubtitle\par}
\end{tcolorbox}

\vspace{1.2em}
\begin{center}
\begin{tabular}{rl}
\textbf{Reviewer:} & \ReviewerName \\
\textbf{Venue:} & \VenueName \\
\textbf{Manuscript ID:} & \ManuscriptId \\
\textbf{Report date:} & \ReportDate \\
\textbf{Recommendation:} & \Recommendation \\
\textbf{Overall score:} & \OverallScore \\
\textbf{Confidence:} & \ReviewConfidence \\
\end{tabular}
\end{center}

\vfill
\begin{summarybox}[Review Scope]
\ReviewScope
\end{summarybox}
\end{titlepage}

\tableofcontents
\newpage

% =============================================================================
\section{摘要与审稿结论}
% =============================================================================

\begin{abstract}
\noindent
TODO: Summarize the manuscript, review scope, evidence base, decisive strengths and weaknesses, final recommendation, and confidence in 1--3 paragraphs.
\end{abstract}

\begin{criticalbox}[Final Recommendation]
\textbf{Recommendation:} \Recommendation

\textbf{Overall score:} \OverallScore

\textbf{Confidence:} \ReviewConfidence

\textbf{Decisive reasons:}
\begin{enumerate}
  \item TODO: decisive reason 1.
  \item TODO: decisive reason 2.
  \item TODO: decisive reason 3.
\end{enumerate}
\end{criticalbox}

% =============================================================================
\section{稿件信息与审查范围}
% =============================================================================

\begin{longtable}{p{0.25\textwidth}p{0.68\textwidth}}
\toprule
\textbf{Item} & \textbf{Content} \\
\midrule
Paper title & \PaperTitle \\
Venue / journal & \VenueName \\
Manuscript ID & \ManuscriptId \\
Reviewed files & TODO: list manuscript, supplement, code, data, rebuttal, or related files. \\
External review backend & TODO: model/backend/thread id, or "not used". \\
Literature search scope & TODO: databases, dates, query families, screening rules. \\
Known limitations & TODO: missing files, inaccessible data, unavailable code, or unverified citations. \\
\bottomrule
\end{longtable}

% =============================================================================
\section{文献检索与创新性审查}
% =============================================================================

\subsection{检索策略与覆盖范围}
TODO: Record database(s), search date, query families, time range, inclusion/exclusion criteria, raw hit counts, deduplication rules, and selected-paper count.

\subsection{最接近工作与差异}
\begin{longtable}{p{0.24\textwidth}p{0.12\textwidth}p{0.25\textwidth}p{0.27\textwidth}}
\toprule
\textbf{Prior work} & \textbf{Year} & \textbf{Overlap} & \textbf{Remaining difference} \\
\midrule
TODO & TODO & TODO & TODO \\
\bottomrule
\end{longtable}

\begin{innovationbox}[Novelty Judgment]
TODO: State whether the contribution is novel, incremental, or substantially overlapped with prior work. Tie every claim to verified evidence.
\end{innovationbox}

% =============================================================================
\section{论文主张与创新点评估}
% =============================================================================

\begin{longtable}{p{0.18\textwidth}p{0.25\textwidth}p{0.22\textwidth}p{0.22\textwidth}}
\toprule
\textbf{Claim} & \textbf{Author evidence} & \textbf{Review assessment} & \textbf{Required fix} \\
\midrule
TODO: C1 & TODO & TODO & TODO \\
TODO: C2 & TODO & TODO & TODO \\
\bottomrule
\end{longtable}

% =============================================================================
\section{方法与逻辑审查}
% =============================================================================

\subsection{整体论证链条}
TODO: Describe the paper's intended reasoning chain from problem to method to evidence to claims.

\subsection{主要逻辑问题}
\begin{criticalbox}[Major Logic Issues]
\IssueHeading{M1}{TODO: issue title}{Major}
\Evidence TODO.

\Impact TODO.

\RequiredAction TODO.

\Confidence TODO.
\end{criticalbox}

% =============================================================================
\section{理论与数学审查}
% =============================================================================

\subsection{定义、符号与目标函数}
TODO: Check notation, definitions, objectives, constraints, and consistency between text and equations.

\subsection{理论声明与证明}
TODO: Check theorem/proposition validity, proof gaps, convergence, complexity, and unsupported theoretical language.

% =============================================================================
\section{实验与实证严谨性审查}
% =============================================================================

\subsection{数据集、划分与评价指标}
TODO: Evaluate datasets, splits, preprocessing, leakage risk, metric choice, and thresholding protocol.

\subsection{基线、消融与统计可靠性}
TODO: Evaluate baseline coverage, head-to-head comparisons, ablations, sensitivity analysis, significance testing, and compute cost.

\begin{evidencebox}[Results-to-Claims Alignment]
TODO: Explain which claims are supported, partially supported, or unsupported by the reported experiments.
\end{evidencebox}

% =============================================================================
\section{写作、图表与可读性审查}
% =============================================================================

\subsection{结构与表达}
TODO: Review paper organization, language clarity, terminology consistency, and missing explanations.

\subsection{图表与复现信息}
TODO: Review figure/table quality, captions, axis labels, reproducibility details, and artifact availability.

% =============================================================================
\section{外部 LLM 审查摘要}
% =============================================================================

TODO: If external review was used, summarize thread id, model/backend, strongest independent criticisms, disagreements, and influence on this report. If not used, state "Not used".

% =============================================================================
\section{综合评价与修改清单}
% =============================================================================

\subsection{核心优势}
\begin{positivebox}[Strengths]
\begin{enumerate}
  \item TODO: strength 1.
  \item TODO: strength 2.
\end{enumerate}
\end{positivebox}

\subsection{主要问题清单}
\begin{longtable}{p{0.08\textwidth}p{0.18\textwidth}p{0.33\textwidth}p{0.30\textwidth}}
\toprule
\textbf{ID} & \textbf{Severity} & \textbf{Issue} & \textbf{Required action} \\
\midrule
M1 & Major & TODO & TODO \\
M2 & Major & TODO & TODO \\
Mo1 & Moderate & TODO & TODO \\
Mi1 & Minor & TODO & TODO \\
\bottomrule
\end{longtable}

\subsection{审稿决议}
TODO: Provide final recommendation, confidence, and minimal conditions for acceptance or resubmission.

\subsection{审稿局限性声明}
TODO: State evidence gaps and what was not evaluated.

\appendix

% =============================================================================
\section{文献检索记录}
% =============================================================================

TODO: Include query strings, raw hit counts, selected references, bibliographic keys, and verification notes.

% =============================================================================
\section{证据矩阵}
% =============================================================================

\begin{longtable}{p{0.10\textwidth}p{0.25\textwidth}p{0.28\textwidth}p{0.25\textwidth}}
\toprule
\textbf{Issue} & \textbf{Evidence source} & \textbf{Interpretation} & \textbf{Required action} \\
\midrule
M1 & TODO & TODO & TODO \\
\bottomrule
\end{longtable}

% =============================================================================
\section{外部审查追踪摘要}
% =============================================================================

TODO: Include trace path(s), thread id(s), and concise summaries or short excerpts within copyright and privacy limits.

\vspace{2em}
\begin{center}
\begin{tcolorbox}[colback=gray!8,colframe=gray!55!black,width=0.82\textwidth,arc=2mm]
\centering
\Large\textbf{审稿报告结束}\\[0.5em]
\normalsize
\textbf{Final recommendation:} \Recommendation\\
\textbf{Overall score:} \OverallScore\\
\textbf{Confidence:} \ReviewConfidence
\end{tcolorbox}
\end{center}

\end{document}
